\section{RQ5: Cross-platform replication}
\label{sec:crossplatform}

The Jetson study above traces the in-app cost to one mechanism: an
accelerator governor that reacts to the duty cycle the host pipeline
creates, coupled to a CPU governor that reacts to how the host waits.
How much of this is specific to that one board? We repeated the core
experiments on two boards with different accelerator architectures
(\cref{tab:boards}): a Raspberry Pi~5 whose Hailo-8 is a fixed-function
INT8 NPU with on-chip NMS and no user-visible frequency control, and a
Rubik Pi~3 (Qualcomm QCM6490), which we measure twice: on its CPU under
a three-cluster \texttt{schedutil} DVFS, and on its Hexagon HTP NPU
through a separately obtained QNN runtime. On both boards we kept the
protocol of
\cref{sec:setup}: official COCOeval on val2017 (all 5,000 images for
accuracy; the ladder uses the first 1,000 on the slower Rubik), the
same letterbox geometry and row conversion, hashed predictions, a
rotated two-repeat campaign, and clock, thermal and (where available)
power logging at 20\,Hz on the Jetson; the Pi's PMIC sustains 10\,Hz
and the Rubik sampler 13--18\,Hz (clocks and thermals only, no power,
\cref{sec:setup}). Two phases deviate from that template, and
we flag them where they appear: the pinned-frequency probes are
single runs per frequency, and the HTP streaming phase assigns the
two rotated repeats one per governor, so each HTP streaming cell is
a single run.

\begin{table*}[t]
\centering
\caption{Replication platforms. The Rubik Pi~3 is measured with two
inference stacks: its CPU, and its Hexagon NPU.}
\label{tab:boards}
\small
\begin{tabular}{@{}llll@{}}
\toprule
 & Raspberry Pi 5 + Hailo-8 & Rubik Pi 3 (CPU) & Rubik Pi 3 (NPU) \\
\midrule
SoC & BCM2712 & QCM6490 & QCM6490 \\
CPU & 4$\times$ A76, one domain, & 8 cores, three domains, & (same) \\
    & 1.5--2.4\,GHz & 0.8--2.7\,GHz & \\
Accelerator & Hailo-8 NPU (INT8), & none used & Hexagon HTP V68 \\
 & on-chip NMS & & (INT8) \\
Model & YOLOv8n HEF & YOLO26n ONNX & YOLO26n QNN ctx, \\
 & & $(1{,}300{,}6)$ & $(1{,}84{,}8400)$, host NMS \\
Runtime & HailoRT 4.23.0 & ONNX Runtime 1.30.0 & qai-appbuilder 2.50.40$^{b}$ \\
CPU governor & \texttt{ondemand}, \texttt{performance} & \multicolumn{2}{l}{\texttt{schedutil}, \texttt{performance}} \\
Power sensor & PMIC rails, SoC side only$^{a}$ & \multicolumn{2}{l}{none usable$^{c}$} \\
\bottomrule
\multicolumn{4}{@{}l@{}}{\scriptsize $^{a}$Excludes the 5\,V loads: the Hailo-8, fan and USB.}\\
\multicolumn{4}{@{}l@{}}{\scriptsize $^{b}$PyPI wheel bundling libQnnHtp + V68 skel; the board's QNN 2.46 predates the}\\
\multicolumn{4}{@{}l@{}}{\scriptsize QAIRT 2.49.40-built context binary.}\\
\multicolumn{4}{@{}l@{}}{\scriptsize $^{c}$The board's battery manager reports inconsistent values; we log clocks and thermals only.}
\end{tabular}
\end{table*}

\subsection{Raspberry Pi 5 with a Hailo-8 accelerator}
\label{sec:pi}

The Hailo model-zoo YOLOv8n HEF offloads the whole network, including
NMS, to the NPU; the host decodes, letterboxes into a uint8 NHWC buffer
and converts the NMS output to COCO rows. The vendor benchmark
(\texttt{hailortcli benchmark}) reports 431.6 images/s in throughput mode,
under both CPU governors; the NPU clock is not user-controllable, and
no register exposes it, so we cannot observe it directly. The per-call
NPU time we can measure includes host-side waiting inside the runtime;
once that share is accounted for, the residual gives no sign of a
duty-cycle response (below). A conventional sequential application on the same HEF
reaches 65.7\,images/s, 15\% of the benchmark. The in-app gap exists on
this accelerator too; with no observable accelerator clock we cannot
rule in or out an internal NPU DVFS, but the mechanism we can identify
is host-side, and the per-call NPU time, once its host-side wait share
is accounted for (below), gives no sign of a duty-cycle response.

\Cref{tab:piladder} shows the pipeline ladder. The mechanisms transfer.
Prefetching decode nearly doubles throughput ($1.86\times$), overlapping
post-processing and adding a second decode worker gives another
$1.5\times$, and a third worker adds nothing, the same saturation point
as on the Jetson. In total the restructuring reaches $2.7\times$ the
sequential application ($3.3\times$ with the spinning wait, see
below), with byte-identical predictions in every rung,
repeat and governor (sha256 prefix \texttt{455ea22b07f1}; official AP
0.320, AP$_S$ 0.122 on all 5,000 images). The wait-policy effect also
appears here, but with a different mechanism than on the Jetson:
replacing the sleeping wait on the asynchronous
inference handle with a spinning wait gains 22\% under
\texttt{ondemand} (177.9 to 217.5\,images/s) and 29\% under
\texttt{performance} (169.9 to 219.8\,images/s), where the CPU clock
is pinned at 2.4\,GHz. The gain therefore cannot be mediated by the
CPU governor; it sits inside HailoRT's wait and callback path, which
we did not profile further. A separate governor effect exists but is
smaller: a controlled probe (\texttt{seq} versus \texttt{seq} with a
busy spinner) shows the governor holding 2.4\,GHz instead of
2.2\,GHz, worth 7.7\% of throughput at 42\% more energy. The NPU time itself, measured from
submission to completion, stays at 6.3--9.3\,ms across the whole ladder
and does not respond to duty cycle; the wait time inside that interval
absorbs the schedule differences. Energy per image falls 1.8$\times$
from sequential to the best variant. We stress that these milliJoules
cover only the SoC rails; the NPU's own draw is on the unmonitored
5\,V input. These values therefore compare SoC-side energy only: they
do not establish whole-system energy rankings, because the unmeasured
Hailo-8 energy can itself move with the schedule. Within-platform
SoC-side comparisons remain meaningful between configurations that share the
same engine and accuracy; cross-engine energy comparisons on this
board need module-level power measurement. Streaming at 15, 30 and
60\,fps is consistent with the Jetson's idle-floor pattern:
sequential and overlapped pipelines differ by 9\% in energy per frame
at 15\,fps (102.0 versus 111.1\,mJ, SoC side) and converge at
60\,fps (37.4 versus 36.2\,mJ). Two cautions apply: the two-repeat
spread on the 15\,fps overlapped cell reaches 19\% (101--121\,mJ),
so that gap is indicative rather than established, and the sequential
pipeline already sits at the idle floor within measurement error (the
idle reading is a single 60\,s window: $1.534$\,W $\Rightarrow$
102.3\,mJ per frame at 15\,fps).

\begin{table}[t]
\centering
\caption{Pipeline ladder on the Pi 5 + Hailo-8 (5,000 val2017 images,
mean of two rotated repeats; p50 is the median per-image end-to-end
pipeline latency; power is SoC-side only).}
\label{tab:piladder}
\footnotesize
\setlength{\tabcolsep}{3.4pt}
\begin{tabular}{@{}llrrrr@{}}
\toprule
Governor & Rung & img/s & p50 ms & SoC mJ/img & CPU MHz \\
\midrule
\multirow{6}{*}{\texttt{ondemand}}
 & seq           &  65.7 & 15.1 & 34.2 & 2020 \\
 & prefetch      & 122.0 & 16.1 & 25.5 & 2243 \\
 & full w1       & 149.8 & 14.6 & 23.0 & 2377 \\
 & full w2       & 177.9 & 25.2 & 25.2 & 2319 \\
 & full w2 spin  & 217.5 & 21.4 & 19.4 & 2300 \\
 & full w3       & 174.7 & 31.5 & 25.2 & 2327 \\
\midrule
\multirow{6}{*}{\texttt{performance}}
 & seq           &  74.6 & 13.3 & 34.8 & 2400 \\
 & prefetch      & 120.5 & 15.1 & 26.1 & 2400 \\
 & full w1       & 140.7 & 14.7 & 22.8 & 2400 \\
 & full w2       & 169.9 & 25.5 & 26.3 & 2400 \\
 & full w2 spin  & 219.8 & 21.2 & 20.0 & 2400 \\
 & full w3       & 166.4 & 32.1 & 27.2 & 2400 \\
\bottomrule
\end{tabular}
\end{table}

Two differences from the Jetson matter. First, accuracy is not
comparable across the boards: the HEF's on-chip NMS is compiled with a
score threshold of 0.2, which truncates the low-score tail that COCO AP
rewards; AP comes out at 0.320, against the 0.373 that Ultralytics
reports for the same network at the conventional 0.001
threshold~\cite{ultralytics2023yolov8}. The threshold explains the
direction and plausibly most of the magnitude, but without pre-threshold
rows we cannot rule out smaller contributions from the HEF toolchain
itself. Because every rung produces
byte-identical predictions, this cancels in all within-platform
comparisons. Second, the default \texttt{ondemand} governor costs less
here than \texttt{nvhost\_podgov} costs on the Jetson: it sits at
2.0\,GHz even in the sequential pipeline and closes most of the gap to
\texttt{performance} once the pipeline is overlapped (219.8 versus
217.5\,images/s for the best rung). The Jetson's threefold governor effect
is specific to its aggressive GPU governor; what generalizes is that
the benchmark number describes a duty cycle the application creates,
not a property of the accelerator. One anomaly we cannot explain from
clocks alone: on the overlap rungs the shipped \texttt{ondemand}
governor is 5--6\% faster than the pinned \texttt{performance}
governor (177.9 versus 169.9\,images/s at \texttt{full\_w2}), while
with the spinning wait the order inverts (217.5 versus 219.8).
NPU-side power states may differ between the two settings; we did not
pursue it.

\subsection{Rubik Pi 3: inference on the CPU}
\label{sec:rubik}

The Rubik Pi~3 runs the same YOLO26n ONNX graph as the Jetson float
engine, but on the CPU (ONNX Runtime, default threads), so inference
itself consumes the resource that the host stages and the wait policy
also need. One structural difference from the accelerator ladders
applies: in the overlap rungs the session call runs in each decode
worker, so up to \texttt{max\_workers} inferences execute concurrently
and the \emph{infer} column includes their mutual contention. The
frequency probes of this subsection pin the run to the prime core
(\texttt{taskset}, intra-op parallelism 1) so that a single frequency
domain is exercised; the ladder runs use the default unpinned
configuration. Accuracy first: on all 5,000 val2017 images the board
reproduces the Jetson engine's AP almost exactly (0.4027 versus
0.40257, AP$_S$ 0.1977 versus 0.19766), so the model is a fixed point
across the two stacks. CPU inference under multithreading emits rows in a nondeterministic
order, so the byte-level hash does not apply directly; canonicalizing
the rows (sorting by image, category, score and box) shows the
predictions are in fact identical: all six rungs produce the same
143,362 canonical rows on the 1,000-image ladder subset in every
repeat and governor (24 of 24 cells), and \texttt{seq},
\texttt{full\_w2} and \texttt{full\_w2\_spin} give identical AP
(0.425).

\Cref{tab:rubikladder} shows the ladder. Every effect that dominated on
the accelerator boards shrinks or inverts. Prefetch alone no longer
helps ($-5$\%): the decode worker steals a core from the inference
threads. Overlap with one worker gives the best result, but only
$+28$\% (versus $3.3\times$ on the Pi and $6.3\times$ on the Jetson,
both under their shipped governors),
and it nearly triples the median latency (304 to 896\,ms)
because the queued frame waits behind a 400\,ms inference. The
spinning wait, worth 20--22\% on both accelerator boards, gains nothing
here: a spinning thread competes with the inference threads for the
same cores. The two governors differ by less than the
thermal drift between identical repeats (3.26 versus 2.97\,images/s for
\texttt{seq}), a picture consistent with both sitting at the thermal
limit: the SoC reached 95.4--95.8\,$^\circ$C --- the firmware ceiling
--- in every ladder run, and a repeat
late in the campaign is up to 17\% slower than an earlier identical
one, a drift the Jetson never showed (median run-to-run
range 0.8\% there). We read no hardware throttle counter on this
board, so we attribute the drift to thermal management by its
signature (sustained ceiling temperature plus progressive slowdown)
rather than by direct observation.

What remains controllable is the clock itself. Pinning the big core's
frequency domain to 806, 1862 and 2707\,MHz moves the sequential
pipeline through 1.7, 3.1 and 3.8\,images/s; inference time scales from
575.7 to 255.2\,ms, a $2.26\times$ ratio for a $3.36\times$ frequency
ratio (one run per pin, and only the prime domain is pinned while the
other two domains float, so we read this as a dose--response rather
than a scaling law). On a CPU-inference device, frequency is a direct
throughput control with no governor mystery: the pipeline never
leaves the processor idle long enough to confuse \texttt{schedutil}. In the streaming
workload (2, 5 and 10\,fps) both pipelines keep up only at 2\,fps,
where the sequential variant again wins on latency (p50 290 versus
356\,ms) and the big core drops to 2.3--2.5\,GHz in the idle gaps. At
5 and 10\,fps both fall behind their arrivals and the queue grows
linearly: within a 300-frame stream the median latency reaches 23.9
and 39.3\,s for the sequential pipeline, while the overlapped
pipeline's higher capacity (3.8 versus 2.8\,images/s) holds it to 8.6
and 26.2\,s, at a permanently saturated 2.7\,GHz. (The streaming
sequential capacity of 2.8\,images/s sits below the ladder's 3.26
because the streaming phase ran late in the campaign, when the board
was hottest.) No schedule rescues
a pipeline whose capacity is below the arrival rate; on this board the
clock, or a cheaper model, is the only lever.

\begin{table}[t]
\centering
\caption{Pipeline ladder on the Rubik Pi 3 (CPU inference, 1,000
val2017 images of which 980 are timed after warm-up, mean of two
rotated repeats with the repeat range in parentheses --- two repeats
cannot support an interval; p50 is the median per-image end-to-end latency,
\emph{infer} the ONNX Runtime call --- in the overlap rungs the decode
workers share the cores that inference uses, so the call itself
lengthens; no power sensor on this board).}
\label{tab:rubikladder}
\footnotesize
\setlength{\tabcolsep}{3pt}
\begin{tabular}{@{}llrrr@{}}
\toprule
Governor & Rung & img/s & p50 ms & infer ms \\
\midrule
\multirow{6}{*}{\texttt{schedutil}}
 & seq           & 3.26 (2.96--3.56) &  304 & 254 \\
 & prefetch      & 3.11 (2.91--3.31) &  641 & 269 \\
 & full w1       & 4.17 (3.99--4.36) &  896 & 415 \\
 & full w2       & 3.85 (3.78--3.92) & 1212 & 446 \\
 & full w2 spin  & 3.84 (3.69--4.00) & 1237 & 451 \\
 & full w3       & 4.04 (3.92--4.17) & 1415 & 426 \\
\midrule
\multirow{6}{*}{\texttt{performance}}
 & seq           & 2.97 (2.89--3.05) &  329 & 275 \\
 & prefetch      & 2.96 (2.86--3.06) &  671 & 281 \\
 & full w1       & 4.07 (3.89--4.26) &  913 & 424 \\
 & full w2       & 4.00 (3.88--4.13) & 1185 & 429 \\
 & full w2 spin  & 3.75 (3.63--3.87) & 1262 & 461 \\
 & full w3       & 3.97 (3.84--4.10) & 1443 & 434 \\
\bottomrule
\end{tabular}
\end{table}

\subsection{Rubik Pi 3 on its Hexagon NPU}
\label{sec:rubikhtp}

The QCM6490's intended accelerator is a Hexagon HTP (V68) NPU, which
the CPU measurements above never touched: the board's QNN runtime
(2.46) predates the available INT8 context binary of the same YOLO26n
graph (built with QAIRT 2.49.40), and the newer SDK requires a vendor
account. We closed the gap with the \texttt{qai-appbuilder} 2.50.40
PyPI wheel, which bundles a matching \texttt{libQnnHtp} and V68 skel;
with the DSP (digital signal processor) library path pointed at the wheel, the context runs
unchanged. Two protocol differences follow from the context, not from
our pipeline: it is INT8 rather than float, and it ends in the
standard $(1{,}84{,}8400)$ head, so NMS runs on the host with the
reference val settings (confidence ${\ge}\,0.001$, IoU $0.7$, at most
300 detections), where the Hailo board sits at the opposite extreme
with NMS fused into the HEF. The quantized stack scores AP $0.364$
(AP$_S$ $0.175$) on all 5,000 images against $0.4027$ ($0.198$) for
the float CPU stack: eleven times the throughput of the best CPU
pipeline under the same governor (43.6 versus 4.07\,images/s under
\texttt{performance}; $6.7\times$ under \texttt{schedutil}) for
$3.9$ points of AP. One protocol caveat applies to that
comparison: the context's uint8 score encoding uses a quantization
scale of $0.01045$ (measured on the output grid), so the effective
score floor is ${\approx}0.01$ rather than the reference $0.001$ --- the same
mechanism we flagged for the Hailo, twenty-fold milder.
It truncates only the far tail (on fewer than ten of the 5,000 images
the stack
emits no rows, all of them images whose float-stack detections score
below $0.02$). A direct sensitivity check quantifies this mechanism alone: quantizing
the float stack's scores to $0.01$ steps and dropping the sub-$0.01$
tail lowers its AP from $0.4026$ to $0.3999$ ($-0.27$ points, AP$_S$
$-0.5$), which is about 7\% of the 3.9-point gap between the two
stacks. The check exercises only the score floor --- not the
quantization of boxes, features or the two stacks' different error
distributions --- so the remaining 93\% of the gap lies elsewhere:
the score floor is one measured contributor, not a bound on the full
INT8 gap. Within the NPU stack, predictions
are byte-identical across every ladder rung, repeat and governor
(sha256 prefix \texttt{09de71d67bad} on the ladder subset), and the
accuracy run's rows for those images match row for row. Each accuracy
row comes from a single run; COCOeval is deterministic, so a repeat
would differ only through runtime nondeterminism.

One deployment finding precedes the measurements: the bundled
appbuilder/QNN userspace wedged the DSP session under concurrent
multi-context inference, so our pipeline holds exactly one context and
issues every NPU call from a single thread; overlap comes from
pipelining decode, input fill and host NMS \emph{around} the
serialized NPU calls. Under that discipline the whole campaign ran
without a single failed call. Separately, a harness bug in an early
version of the evaluation script inflated every timed call by a
per-call error print; we fixed it, reran the entire HTP campaign, and
report only the corrected numbers. Both issues, and why the bug
mattered most exactly where the paper looks at clock sensitivity, are
detailed in \cref{app:qnn}.

\Cref{tab:rubikhtpladder} shows the ladder, and the CPU board's
picture inverts back to the accelerator pattern. With a
24--28\,ms NPU call, the host stages are again the tax, and the
ladder recovers them: $+24$\% under \texttt{schedutil} (22.9 to
28.5\,images/s, the spinning-wait diagnostic; the best
non-spinning rung reaches 27.8) and $+52$\% under
\texttt{performance} (28.8 to
43.6\,images/s), where the best rung saturates the accelerator: its
22.9\,ms period matches the effective per-call service time under
load, marginally below the 24.4\,ms isolated-call median. One
consequence is that a second decode worker does not pay off under
\texttt{performance}: \texttt{full\_w2} is slower than
\texttt{full\_w1} (39.9 versus 43.6\,images/s) and adds 39\,ms of p50
latency, because NPU calls are serialized on one thread, so the extra
worker cannot add accelerator parallelism and only lengthens the
queue; under \texttt{schedutil} the two rungs tie (27.8 versus
27.7\,images/s), the slowed feeder core being the binding constraint
either way. That
median corresponds to about 41 calls/s of back-to-back submission,
and the best rung exceeds it (43.6\,images/s), so the isolated call
overstates the loaded service time --- under sustained submission
either queued calls overlap driver-side work or the loaded dispatch
share shrinks; we cannot separate the two, and report both the
isolated call and the saturated period. The governor,
irrelevant on the CPU board, is here the largest single term:
under \texttt{schedutil} the core that feeds the NPU settles at
$0.95$--$1.0$\,GHz during the overlapped rungs and at $1.78$\,GHz in
the sequential one --- the less host work a rung leaves on the
critical path, the lower the clock the governor chooses --- and the
same pipeline runs up to 57\% faster under \texttt{performance}
(43.6 versus 27.8\,images/s on the matched \texttt{full\_w1} rung;
44\% on the matched \texttt{full\_w2} rung, 39.9 versus
27.7\,images/s).
The pinned-frequency probe shows why the host clock matters so much
for an ``NPU'' call: at 806, 1862 and 2707\,MHz the sequential
pipeline delivers 16.7, 24.0 and 28.6\,images/s, and the measured NPU
time itself falls from 35.8 to 24.4\,ms. The three points are
consistent with a decomposition into a fixed ${\approx}20$\,ms NPU
kernel plus a host-side dispatch term inversely proportional to the
pinned core's clock (least-squares fit; the mid point sits
$0.8$\,ms above it, so the two-term form is descriptive rather than
exact): a large share of each
accelerator call is remote-procedure-call (RPC) work executed on the CPU. We stress that this
is a descriptive decomposition, not a validated model --- two
parameters over three single-run pins leave one residual degree of
freedom, and because the NPU exposes no observable clock or load
register, the decomposition cannot distinguish a clock-invariant
kernel from an accelerator DVFS whose response happens to be absorbed
by the constant term. The governor invariance we report below is a
hypothesis consistent with the data, not a measurement. The spinning
wait gains nothing on this stack (38.2 versus
39.9\,images/s under \texttt{performance}; it also shows the widest
repeat spread of the campaign, 34.9--41.5\,images/s across the two
repeats --- thermal drift, with the SoC climbing from 83 to
95\,$^\circ$C, not noise): with a single serial caller there is no
asynchronous-handle wait for the spinner to heat. As on the Hailo,
the NPU exposes no user-visible frequency control, and within this
decomposition the accelerator clock is not the moving part: what moves
with the governor is the call's host-dispatch share. The sequential call
measures 24.4\,ms under \texttt{performance} and 28.3\,ms under
\texttt{schedutil}, while the pinned probe's kernel component
(${\approx}20$\,ms) is governor-invariant in the fit; correspondingly, under \texttt{performance} the best
rung issues one call every 22.9\,ms, the loaded per-call service time,
while under the shipped governor the best rung issues one every
35.1\,ms against a 28.3\,ms sequential call, so the accelerator idles
about a fifth of each period even when the queue never empties. (The \emph{infer} column of
\Cref{tab:rubikhtpladder} grows from 24.4 to ${\approx}47$\,ms in the
overlap rungs because there the measurement spans the queue wait
behind the single serialized NPU caller, not because the accelerator
slows down.) Streaming behaves as the capacities predict. Under
\texttt{schedutil} the pipeline keeps up at 15\,fps (p50 108\,ms
overlapped, 54\,ms sequential) and queues linearly at 30 and 60\,fps
(p50 0.8 and 5.1\,s overlapped, 2.4 and 6.7\,s sequential, within a
480-frame stream). Under \texttt{performance} the overlapped pipeline
sustains 30\,fps at p50 47\,ms while the sequential one (capacity
28.8\,images/s) already falls behind (p50 0.39\,s); at 60\,fps both queue
(p50 2.1 versus 4.7\,s). Predictions agree row for row with the
ladder on the shared images, in both campaigns. Each streaming cell here is a single run (the two
rotated repeats were assigned one per governor), so the governor
contrast is paired but not replicated.

\paragraph{Baking selection into the graph}
The Jetson and the Rubik CPU stacks run the NMS-free $(1{,}300{,}6)$
head with selection fused into the model (the Hailo HEF fuses its own
NMS on-chip instead), so for completeness we also
built an end-to-end context for the HTP with the matching QAIRT
2.50.40 toolchain (INT8, 500 calibration images, the context binary
generated on the board itself). One silent pitfall is worth
recording: the exported head routes the box branch
(range ${\approx}\,[-50,800]$) and the score branch ($[0,1]$) through
a single concat--transpose--split, and per-tensor INT8 quantization
gives the shared tensor one scale, which sets the score channel to
zero at runtime. Calibration cannot fix this; graph surgery can ---
we transpose the two branches separately and expose boxes, scores and
classes as three output tensors, so each receives its own scale (the
V68 HTP rejects fp16 graphs outright, so mixed precision was not an
option). The repaired context detects correctly and is faster end to
end, because selection moves onto the DSP and the host NMS stage
disappears: the same descriptive decomposition puts the fixed NPU
kernel at ${\approx}11.6$\,ms (endpoint and least-squares fits
span 11.5--11.6\,ms, with the mid probe point $0.3$\,ms above the
fit) against ${\approx}20$\,ms for the standard
head, the accuracy run
rises to 32.5\,images/s and the
best rung to 63.2\,images/s (+45\%), and the overlapped 30\,fps stream
cuts its p50 latency by 42\% under \texttt{performance} (27 versus
47\,ms). It also costs accuracy: AP $0.293$ against $0.364$, with the
loss concentrated in small objects (AP$_S$ $0.112$ versus $0.175$),
where the in-graph INT8 box-decode chain (regression logits at
$0.75$ per step, decoded boxes at $3.3$\,px per step) hurts most
(\Cref{tab:rubikhtpe2e}). We therefore keep the standard head with
host NMS as the primary HTP stack and report the fused head as a
measured configuration comparison: on this toolchain, the
fused-selection configuration runs 45\% faster at peak and scores 7
AP points lower than the standard-head configuration. The comparison
is not a clean ablation --- the two contexts come
from different QAIRT releases (2.49.40 versus 2.50.40) and different
calibration sets --- so the $-7$ AP is the observed difference between
these two configurations, not an isolated cost of fusing selection.

\begin{table}[t]
\centering
\caption{Pipeline ladder on the Rubik Pi 3 with the Hexagon HTP NPU
(INT8 context, host NMS; 1,000 val2017 images of which 980 are timed,
mean of two rotated repeats with the repeat range in parentheses ---
two repeats cannot support an interval; NPU calls serialized on one
thread, see text). \emph{infer} is submission-to-completion of the
NPU call; in the overlap rungs it spans the queue wait behind the
single serialized caller, not pure accelerator time (the sequential
rung gives the pure call: 24.4\,ms under \texttt{performance}).}
\label{tab:rubikhtpladder}
\footnotesize
\setlength{\tabcolsep}{3pt}
\begin{tabular}{@{}llrrr@{}}
\toprule
Governor & Rung & img/s & p50 ms & infer ms \\
\midrule
\multirow{6}{*}{\texttt{schedutil}}
 & seq           & 22.94 (22.82--23.06) &  43 & 28.3 \\
 & prefetch      & 25.60 (25.52--25.68) &  77 & 28.8 \\
 & full w1       & 27.75 (27.75--27.76) & 133 & 47.0 \\
 & full w2       & 27.68 (27.65--27.71) & 173 & 46.2 \\
 & full w2 spin  & 28.45 (28.26--28.64) & 173 & 53.2 \\
 & full w3       & 27.70 (27.63--27.77) & 209 & 46.7 \\
\midrule
\multirow{6}{*}{\texttt{performance}}
 & seq           & 28.75 (28.67--28.83) &  35 & 24.4 \\
 & prefetch      & 31.86 (31.74--31.98) &  62 & 24.4 \\
 & full w1       & 43.59 (43.56--43.61) &  89 & 40.0 \\
 & full w2       & 39.89 (39.54--40.24) & 128 & 45.3 \\
 & full w2 spin  & 38.22 (34.92--41.53) & 135 & 47.2 \\
 & full w3       & 40.01 (37.47--42.56) & 146 & 42.5 \\
\bottomrule
\end{tabular}
\end{table}

\begin{table*}[t]
\centering
\caption{Head configuration comparison on the Rubik HTP: the standard
head with host NMS versus the NMS-free head fused into the INT8
context. This is not a clean ablation --- the two contexts come from
different QAIRT releases (2.49.40 versus 2.50.40) and different
calibration sets --- so the accuracy difference is a property of the
two configurations, not an isolated cost of fusing selection.
Accuracy on all 5,000 val2017 images (same pipeline, serialized NPU
calls); throughput is the accuracy run and the best ladder rung per
governor. Under \texttt{schedutil} both best cells are the
spinning-wait diagnostic (best non-spinning rungs: 27.8 and 36.9
images/s); under \texttt{performance} the fused head's spin and
non-spin best rungs tie within run-to-run noise (63.2 versus 63.0).}
\label{tab:rubikhtpe2e}
\small
\begin{tabular}{@{}lccccc@{}}
\toprule
Head & AP & AP$_{50}$ & AP$_S$ & acc.\ img/s & best img/s (sched/perf) \\
\midrule
standard $(1{,}84{,}8400)$, host NMS & 0.364 & 0.527 & 0.175 & 22.8 & 28.5 / 43.6 \\
NMS-free, fused (3 outputs)        & 0.293 & 0.477 & 0.112 & 32.5 & 39.6 / 63.2 \\
\bottomrule
\end{tabular}
\end{table*}

\subsection{What transfers, and what does not}
\label{sec:transfer}

Across three boards and four inference stacks the same experiment
gives four shades of one answer. The \emph{benchmark-versus-application
gap} appears on every stack we measured: 431.6 images/s versus
65.7\,images/s on the Pi~5 (15\%), 4.2\,ms versus 13.2\,ms on the
Jetson, and on the Rubik a 254\,ms inference inside a 304\,ms
sequential loop on the CPU and, under the shipped governor, a 28\,ms
NPU call inside a 43\,ms loop on the HTP. These four gaps are not
all of the same kind: each compares an isolated engine or
inference-call benchmark with a sequential application, but the content
of the difference is stack-specific. On the Jetson it is mostly the
governor's response to the application's duty cycle
(\cref{sec:diagnosis}); on the Pi, where the accelerator call is
fixed-cost, it is the host work around that call; on the HTP it
additionally contains the host-dispatch share of each accelerator
call and that share's response to the CPU governor; on the Rubik CPU
it is bounded by thermal saturation. The \emph{host-pipeline ladder} helps on all four stacks,
but its size depends on what else moves with the schedule (best rung
versus sequential, each under the board's shipped governor):
$6.3\times$ on the Jetson against the naive sequential rung (float
input, stream synchronization; $4.3$--$5.3\times$ against the stronger
sequential variants of \cref{tab:ladder}; the $2.4\times$ headline of
\cref{sec:intro} uses the locked-clock baseline), where overlapping
also lifts the GPU out
of its DVFS floor, so the ladder undoes both the host-stage tax and
the clock collapse; $3.3\times$ on the Pi, whose sequential pipeline is
the most host-bound of the four; $1.3\times$ on the Rubik CPU; and
only $1.2\times$ on the HTP under the shipped governor, where the
CPU governor settles the core that feeds the DSP near 1\,GHz as soon
as the pipeline overlaps --- fixing the governor first re-opens the
ladder to $1.5\times$ (\texttt{performance} column of
\Cref{tab:rubikhtpladder}). The \emph{wait-policy effect} appears on
both asynchronous-handle accelerators but at different rungs and through
different mechanisms. On the Jetson it is a property of the schedules
that actually wait: replacing the sleeping wait gains 20--25\% on the
sequential and prefetch rungs, mediated by the CPU governor (the
busy-loop and \texttt{uclamp} experiments of \cref{sec:sched}), and
nothing on the full two-worker pipeline, whose host never goes idle on
the critical path. On the Pi the same swap gains 22--29\% on the
\emph{best} rung (\texttt{full\_w2}), and it persists at a pinned CPU
clock, so it lives in the runtime's wait path rather than the
governor. It vanishes where no such wait exists: when
the CPU itself computes, because the spinner steals cycles from the
inference it waits for, and on the HTP stack, where the single serial
caller waits on a bounded queue rather than an async handle. The
\emph{accelerator governor}, the Jetson's dominant term, is observable
only where the accelerator exposes a clock: neither NPU offers a
user-visible clock or load register, so an internal accelerator DVFS
cannot be ruled out --- but it is not needed to explain these
measurements: the per-call NPU time we can measure includes the
host-side wait and queueing inside the runtime (it grows from 6.3 to
9.3\,ms across the Pi's ladder, and the HTP's infer column spans queue
waits), and once that share is accounted for, neither stack shows a
residual that moves with the host duty cycle
--- on the HTP the overlap rungs saturate the loaded per-call service
time under \texttt{performance} (under the shipped
\texttt{schedutil}, whose slowed core stretches the call's
dispatch share, the accelerator still idles about a fifth of each
period at the best rung; the \emph{infer} column of the
ladder tables grows in the overlap rungs because it spans queueing
behind the single serialized caller, not because the accelerator
slows down; \cref{sec:rubikhtp}). But the governor story does not
disappear on the NPU boards, it moves: on the HTP stack the
\emph{CPU} governor becomes the largest term (57\% on the matched
\texttt{full\_w1} rungs), because
\texttt{schedutil} reads the idleness the accelerator creates and
down-clocks the very core that feeds the DSP, and every ``NPU'' call
carries a host dispatch share that grows as the core clock falls. Finally, sustained
CPU inference introduces a further constraint: the SoC sits at its
95\,$^\circ$C ceiling and otherwise identical repeats separate
measurably. The NPU stacks approach the same ceiling under
\texttt{performance} (91--95\,$^\circ$C). The Pi shows no drift, but on
the HTP the two heaviest \texttt{performance}-governor overlap rungs
do: \texttt{full\_w2\_spin} falls from 41.5 to 34.9\,images/s between
repeats as the SoC climbs from 83 to 95\,$^\circ$C, and
\texttt{full\_w3} from 42.6 to 37.5, so those two cells mix the
governor effect with a thermal one (the \texttt{schedutil} rows and
both sequential ladders show no drift). On the CPU stack the thermal limit is the binding constraint. And \emph{streaming behavior} follows
capacity on the NPU boards exactly as \cref{sec:stream} predicts
(\Cref{tab:streamreplication}): the Pi keeps up at every rate under
either pipeline --- its sequential capacity of 65.7\,images/s already
exceeds 60\,fps --- while the HTP queues precisely when the arrival
rate passes the capacity the ladder measured. One latency mechanism
does not transfer literally. The Jetson's plain overlap holds each
frame's rows until the next frame is enqueued, adding about one
arrival period (\cref{sec:stream}); that hold is a property of our
double-buffered host loop, and it is what the arrival-aware variant
removes. Neither NPU stack has it: HailoRT delivers each frame's
results through an asynchronous completion callback, so the Pi's
overlapped p50 tracks the service time (${\approx}17$\,ms) at every
rate, and the HTP pipeline submits rows as soon as the serialized
call returns. There the shipped governor is what inflates latency:
the slowed feeder core stretches both the call and the host NMS that
overlaps it, so at 15\,fps the overlapped p50 is 108\,ms under
\texttt{schedutil} against 42.6\,ms under \texttt{performance}. What
transfers unchanged is the capacity-margin behavior: queues build
exactly when the arrival rate passes the ladder-measured capacity.
For a practitioner the
ordering of
\cref{sec:discussion} survives with two amendments: on CPU-inference
devices, stop after decode overlap with a single worker and spend the
effort on clocks and cooling; on NPU devices, check the CPU governor
early --- on the QNN/HTP stack it was the largest single term, while
the Hailo stack's shipped governor was already within 1\% of the
pinned one at the best rung.

\textbf{Answer to RQ5.}
The findings carry over in outline, not in mechanism. Across all four
stacks (Jetson TensorRT, Hailo-8, Rubik CPU, Rubik HTP), isolated
engine or inference-call measurements understate the end-to-end
application cost, but for different reasons: on the Jetson the gap is
a GPU-governor-induced engine slowdown; on the Pi it is host and
runtime overhead around a fixed-function NPU; on the Rubik CPU it is
shared CPU resources and thermal saturation; and on the HTP it is the
CPU-side dispatch share of each call and that share's response to the
CPU governor. In each case the host-pipeline ladder recovers
throughput without touching the engine, and streaming queues exactly
when the arrival rate passes the ladder-measured capacity. Across model architectures
(\cref{sec:generality}), YOLOv10n and YOLO26s (whose engine takes
$1.6\times$ longer per image)
reproduce the same schedule effects ($5.6$--$6.9\times$ over the
Ultralytics predictor), so the conclusions are not specific to YOLO26n.

\begin{table*}[t]
\centering
\caption{Streaming on the NPU boards: p50 latency per frame in ms
(release to finished rows; 1,000-frame streams on the Pi, 480-frame
streams on the HTP). Pi: mean of two rotated
repeats under the shipped \texttt{ondemand} governor, with SoC-side
energy per frame in parentheses. HTP: one run per cell and governor
(the two rotated repeats were assigned one per governor), so the
governor contrast is paired but not replicated; the fused-head variant
is discussed in the text. Queueing cells are
rounded to two significant figures.}
\label{tab:streamreplication}
\small
\begin{tabular}{@{}lllrrr@{}}
\toprule
Board & Pipeline & Governor & 15 fps & 30 fps & 60 fps \\
\midrule
Pi 5 + Hailo-8 & seq     & ondemand    & 17.4 (102.0\,mJ) & 16.9 (61.0\,mJ) & 15.7 (37.4\,mJ) \\
               & full w2 & ondemand    & 17.3 (111.1\,mJ) & 17.0 (58.7\,mJ) & 15.6 (36.2\,mJ) \\
\midrule
Rubik HTP      & seq     & schedutil   &  54.2 & 2400 & 6700 \\
               &         & performance &  35.2 & 390 & 4700 \\
               & full w2 & schedutil   & 108 & 820 & 5100 \\
               &         & performance & 42.6 & 47 & 2100 \\
\bottomrule
\end{tabular}
\end{table*}
